\section{Introduction}

A diluted spin glass has a finite mean number of interactions per spin.
Its local environment remains random as the system grows, and the cavity
method describes that environment through a distribution of local fields.
When many competing states contribute to the partition function, the
proposed order parameter is a hierarchy of distributions of these fields.
The thermodynamic question is whether this hierarchy determines the
limiting expected log partition function per spin. We call this quantity
\emph{pressure}. At inverse temperature $\beta$, the usual physical free
energy per spin is its negative divided by $\beta$.

The Viana--Bray model is a classical diluted two-spin
model~\cite{VB85}. M\'ezard and Parisi~\cite{MP01} developed the
hierarchical cavity picture for the fixed-connectivity Bethe-lattice
spin glass, with an explicit one-step replica-symmetry-breaking
functional. Franz and Leone~\cite{FL03} adapted Guerra's
interpolation method~\cite{Guerra03} to diluted systems and obtained
replica-symmetric and one-step pressure bounds for even-arity models.
Panchenko and
Talagrand~\cite{PT04} then proved the arbitrary finite-level
hierarchical upper bounds considered here.

We prove that the infimum of these bounds is the limiting pressure.
This resolves the equality conjecture stated immediately after
Theorem~4 of Panchenko--Talagrand~\cite{PT04}, for their Poisson,
even-arity Ising model class with factorized interactions and the
specified positivity condition. Only first moments of the interaction
and external field are required. The result includes the Viana--Bray
model, diluted even-spin models with symmetric couplings, and soft
random even-$K$ SAT at every positive temperature. The exact model and
trial functional are defined in \cref{model:section}; the named models
and their expected ground-state limits are treated in
\cref{models:section}.

An exact variational description and reduction to a prescribed
trial class are distinct problems. The cavity variational principle of
Aizenman, Sims and Starr~\cite{ASS03} for the Sherrington--Kirkpatrick model
has a Viana--Bray counterpart in De Sanctis's formulation over random
multioverlap structures~\cite{DeSanctis04}. Panchenko's
spin-distribution approach~\cite[Theorems~1--2]{Pan13} gives an exact
variational formula over invariant exchangeable spin distributions.
Under their positivity condition, Coja-Oghlan and
Perkins~\cite{COP19} obtain an invariant-kernel variational formula
for random regular factor graphs. These results
use different descriptions of the state from the finite hierarchy of
message laws optimized here. The additional task is to recover the
specific hierarchical cavity functional without losing the pressure.

Panchenko~\cite[Theorem~2]{Pan16} established a hierarchical
single-site representation for finite-RSB asymptotic Gibbs measures of
suitably modified models; approximation of general states by that class
remains a separate step. Biswas, Chen and
Sen~\cite[Theorem~1.1 and Proposition~1.4]{BCS25} obtain exact
replica-symmetric pressure formulas in high-temperature and critical or
subcritical regimes for compact real spin spaces under moment and regularity
assumptions. For the zero-field Ising spin glass with Rademacher couplings
on random regular graphs, Concetti~\cite[Theorem~1.3]{Concetti25}
gives a full-replica-symmetry-breaking pressure bound no larger than each
optimized finite-level bound; equality with that functional is
conjectural in that work. The result here
identifies the finite-depth Panchenko--Talagrand infimum at arbitrary
density in its stated interaction class.

The main obstacle is the distribution of spins seen across several
replicas, that is, several samples from the same Gibbs system. The
probability of each replica-spin pattern depends on products of every
subset of those spins. The spatial average of such a product is a
\emph{multioverlap}; pair overlaps alone do not determine all these
patterns. This is the information isolated by the spin-distribution
approach~\cite{Pan13}. To compare a cavity insertion with independent
hierarchical messages, we must control all finite multioverlaps before
redrawing the message roots independently.

We work with finite auxiliary hierarchies before taking the
thermodynamic limit. Their branching is part of the construction;
no ultrametric organization of the physical Gibbs states is assumed.
Bounded perturbations yield identities for
prescribed replica branching patterns. Delaying each branching depth
by one level permits an induction that concentrates multioverlaps
\emph{on average over branching depths}. The cavity insertion itself
has an expansion averaging over finite trees, so this averaged control
is sufficient. The system size tends to infinity at each fixed depth;
only then does the depth increase. The proof strategy following
\cref{main:theorem} explains how the finite-tree calculus, concentration
and independent-message comparison fit together.
